\chapter{Bronchoscopy}

\section*{What Is This Test or Procedure?}
Bronchoscopy is a procedure that examines the airways (trachea and bronchi) with a thin, flexible tube passed through the nose or mouth into the lungs.

\section*{Alternative Names}
Flexible bronchoscopy; rigid bronchoscopy; bronchoscopic biopsy

\section*{Principle}
A bronchoscope with a camera is advanced into the airways to visualize the mucosa. Instruments passed through the working channel allow sampling (biopsy, brushing, bronchoalveolar lavage), aspiration, stent placement, or removal of foreign bodies.

\section*{History}
Rigid bronchoscopy was introduced in the 1890s by Gustav Killian, and the flexible fiberoptic bronchoscope was developed by Shigeto Ikeda in the 1960s, transforming pulmonary medicine.

\section*{Why This Test or Procedure Is Ordered}
Performed to evaluate persistent cough, hemoptysis, abnormal imaging findings, suspected infection or malignancy, and to obtain biopsy specimens, clear mucus plugs, or remove aspirated foreign bodies.

\section*{Is This a Primary or Secondary Test or Procedure}
Primary diagnostic and therapeutic procedure for airway and lung disease.

\section*{How This Test or Procedure Is Performed}
Performed under topical anesthesia with sedation, or general anesthesia for rigid bronchoscopy. The scope is passed through the nose or mouth into the trachea and bronchi; samples are taken as needed. The patient is monitored throughout.

\section*{What Preparation Is Needed}
NPO before sedation, review of anticoagulants, and pre-procedure imaging. Patients with significant respiratory compromise may require extra precautions.

\section*{Contraindications and Precautions}
Severe refractory hypoxemia, uncontrolled bleeding, recent myocardial infarction, and unstable arrhythmia are relative contraindications. Precaution: bleeding risk increases with biopsy in uremia or coagulopathy.

\section*{Risks}
Bleeding, pneumothorax, hypoxemia, fever, laryngospasm, and transient sore throat.

\section*{What Happens After the Test or Procedure}
The patient is observed until sedation wears off and is asked to fast until the gag reflex returns. Specimens are sent for culture, cytology, and histology.

\section*{Understanding Your Results}
Cytology or biopsy may reveal malignancy, infection, inflammation, or granulomas; lavage culture identifies pathogens.

\section*{Normal Results}
Patent airways with intact mucosa and no abnormal masses, secretions, or foreign bodies.

\section*{Follow-up Tests and Procedures}
Results guide treatment; repeat bronchoscopy may be needed for surveillance or therapy.

\section*{References}
\begin{enumerate}
  \item MedlinePlus. Bronchoscopy. U.S. National Library of Medicine; 2025.
  \item Current Medical Diagnosis and Treatment. McGraw-Hill; 2025.
\end{enumerate}

\clearpage
